\chapter{Conclusiones}

Habiendo implementado el circuito propuesto, se verificaron experimentalmente las características intrínsecas de los
amplificadores con retroalimentación negativa frente a su modelo simulado y teorico. Los niveles de ganancia medidos presentaron
una alta concordancia con los valores teóricos, registrando un error porcentual de apenas 9.15\% respecto a la
simulación, y evidenciando la caída de ganancia en lazo cerrado característica de esta topología.

En cuanto a las impedancias, estas también reflejaron las mejoras proyectadas, midiendo -.

Por otro lado, la medición de desensibilidad demostró empíricamente la estabilidad del diseño: se obtuvo una variación
de solo el 13\% en la ganancia de lazo cerrado frente a una drástica variación del 55\% en lazo abierto, confirmando la
estabilidad del sistema.

Finalmente, el análisis de la respuesta en frecuencia arrojó resultados precisos, visualizándose la atenuación de la
ganancia acompañada por el consecuente aumento en el ancho de banda, validando así los compromisos fundamentales en el
diseño de esta clase de amplificadores.

\vspace{1cm}
\noindent\textbf{Principales observaciones:}
\begin{itemize}
    \item [OBSERVACIÓN 1]
    \item [OBSERVACIÓN 2]
    \item [OBSERVACIÓN 3]
\end{itemize}
